
Reinforcement learning from human feedback (RLHF) alignment is known to degrade the calibration of large language models, yet the underlying mechanism remains unresolved. This study presents a pre-registered mechanistic discrimination among three candidate hypotheses: monotonic scale inflation (H1), decision-boundary restructuring (H2), and framing susceptibility (H3). Experiments were conducted on the Pythia alignment ladder (1.4B, 2.8B, and 6.9B parameters) with SFT, DPO, and PPO variants, evaluated on MMLU (14,042 items) and TruthfulQA MC1 (817 items) using lm-eval-harness v0.4.11. The principal finding is that H1 is not supported within this model family: all nine Spearman rank correlations between base and aligned log-probability distributions fall below the pre-registered H1 threshold of rho >= 0.90, and eight of nine fall below 0.85. Under PPO alignment at 1.4B parameters, rho = -0.324 and 99.7\% of MMLU items receive a different top-ranked answer after alignment. H3 is not supported in the softmax-ECE evaluation setting, as MMLU calibration degradation exceeds TruthfulQA calibration degradation for all alignment types. These results are consistent with H2 (decision-boundary restructuring) as the dominant mechanism within the Pythia 1.4B-6.9B family. An exploratory observation indicates that DPO produces larger calibration degradation than PPO across all three model sizes, contrary to the pre-registered ordering prediction; this finding is based on public fallback checkpoints and requires replication with matched-training checkpoints. One planned sub-hypothesis (H-M4, testing ATS correction effectiveness) was not executed.
